\documentclass[11pt,a4paper]{article}

\usepackage[T1]{fontenc}
\usepackage[utf8]{inputenc}
\usepackage{natbib}
\usepackage{lmodern}
\usepackage{geometry}
\usepackage{hyperref}

\geometry{margin=2.5cm}
\hypersetup{
  colorlinks=true,
  linkcolor=black,
  citecolor=black,
  urlcolor=blue,
  pdftitle={State of the Art on Synthetic Populations and LLMs},
  pdfauthor={Codex}
}

\title{State of the Art on Synthetic Populations and LLMs}
\author{Academic synthesis document}
\date{March 9, 2026}

\begin{document}

\maketitle

\begin{abstract}
Synthetic populations are a core component of microsimulation, transportation modeling, public policy analysis, and agent-based simulation. Their purpose is to reconstruct an artificial population whose statistical properties are consistent with aggregate constraints, individual samples, or heterogeneous data sources. In parallel, large language models (LLMs) open a new methodological direction: not only generating plausible sociodemographic structures, but also simulating behaviors, opinions, and interactions conditioned on individual profiles. This state-of-the-art review distinguishes classical approaches to population synthesis, examines the contribution of LLMs as behavioral engines, and discusses the conditions under which these two traditions can be combined. The analysis shows that LLMs do not replace statistical reconstruction methods; at best, they complement them by modeling latent, narrative, and decision-oriented dimensions, while introducing substantial risks in terms of bias, calibration, and external validity.
\end{abstract}

\textbf{Keywords:} synthetic population, microsimulation, agent-based simulation, agent-based modeling, large language models, silicon samples.

\section{Introduction}
The concept of a synthetic population refers to a set of artificial individuals and, often, households constructed in such a way that they reproduce the characteristics of a real population from incomplete, partial, or anonymized data. The goal is not to replicate actual persons, but to generate a statistically coherent representation of a target population that can be used in models of mobility, consumption, public health, urban planning, or epidemic spread \citep{beckman1996,chapuis2022}.

Historically, the literature on synthetic populations has focused on three major questions: how to satisfy aggregate margins, how to reconstruct credible multivariate dependencies, and how to guarantee structural coherence across individuals, households, and spatial units. Since the early 2020s, a second body of literature has emerged around LLMs, which are used to simulate synthetic respondents, personas, or more autonomous agents in multi-agent environments \citep{argyle2023,park2023,gao2024}. This development shifts the problem: it is no longer only a matter of whether a population statistically resembles reality, but also whether it can be endowed with plausible behaviors in tasks involving interaction, decision-making, or survey response.

This document argues that these two traditions are complementary rather than interchangeable. Synthetic populations primarily belong to a framework of constrained reconstruction and calibration. LLMs, by contrast, operate more naturally as generators of conditional content, useful for enriching already structured agents, yet still fragile when representativeness, calibration, and behavioral stability are required.

\section{Foundations and Objectives of Synthetic Populations}
In most applications, the construction of a synthetic population serves at least one of the following objectives:
\begin{enumerate}
  \item reproducing observable marginal and joint distributions for a given geographical area;
  \item creating microdata suitable for microsimulation or agent-based models;
  \item protecting confidentiality by replacing sensitive microdata with artificial records;
  \item enabling counterfactual analysis when exhaustive individual-level data are unavailable.
\end{enumerate}

The quality of a synthetic population therefore cannot be reduced to statistical plausibility alone. In advanced applications, one must also consider the coherence of within-household relations, spatial consistency, stability across scales, and the capacity of the resulting population to feed downstream simulators without introducing major artifacts \citep{farooq2013,chapuis2022}.

\section{Classical Approaches to Population Synthesis}
\subsection{Reweighting and IPF-based methods}
One of the most influential traditions starts from an individual sample, for example a public-use microdata sample or census extract, and then reweights or selects individuals so as to satisfy local aggregate constraints. The seminal work of \citet{beckman1996} formalized this framework in the context of activity-based and transportation models. Iterative Proportional Fitting (IPF), together with its multivariate extensions, plays a central role here: contingency tables are adjusted iteratively in order to satisfy known margins.

These approaches are robust, interpretable, and efficient whenever a sufficiently rich reference sample is available. Their main limitation is equally well known: they inherit the weaknesses of the source sample, including empty cells, under-representation of rare combinations, and limited capacity to reconstruct higher-order dependencies.

\subsection{Combinatorial and simulation-based approaches}
A second family of methods formulates synthesis as a selection or combinatorial optimization problem. Instead of merely adjusting weights, the objective is to construct an entire synthetic micro-population that minimizes the discrepancy with observed constraints. In this spirit, \citet{farooq2013} propose a simulation-based approach supported by Markov chain Monte Carlo mechanisms. The appeal of such methods is twofold: they allow a broader exploration of the solution space and make it possible to generate several plausible populations rather than a single reconstruction.

Compared with pure IPF, these methods often handle incompatibilities among constraints more effectively, but they are also more computationally costly, more sensitive to parameter settings, and generally harder to calibrate at large scale.

\subsection{Sample-free approaches}
When individual microdata are unavailable, the literature has developed so-called \emph{sample-free} approaches. \citet{lenormand2012} provide an early systematic comparison of sample-based and sample-free methods for generating populations structured into households. \citet{ye2017} extend this line of work by inferring joint distributions from margins and partial distributions, together with a theoretical discussion of multivariate IPF convergence.

The practical relevance of such approaches is considerable in institutional settings where microdata cannot be disseminated. Their main weakness lies in identifiability: identical aggregate constraints may be compatible with multiple distinct microstructures. In other words, the resulting solution may be plausible without being strongly determined.

\subsection{Recent reviews and opening toward learning-based methods}
The review by \citet{chapuis2022} shows that the field has progressively moved beyond a simple opposition between sample-based and sample-free methods. Recent work seeks to integrate multiple data sources, multiple geographical scales, and multiple entities such as individuals, households, dwellings, and activities. There is also a growing interest in learned generative models for reconstructing complex dependencies or continuous attributes, even though such approaches remain less standardized than IPF and combinatorial optimization in many public-sector applications.

\section{How to Evaluate a Synthetic Population}
Evaluation remains a recurrent source of tension in the field. A synthetic population may satisfy observable margins perfectly while still distorting the latent structures that matter for downstream simulation. In practice, four levels of evaluation are common.

First, distributional validity measures the discrepancy with observed margins and cross-tabulations. Second, structural validity concerns the coherence of links among individuals, households, and dwellings. Third, functional validity examines the effects of the synthetic population on the outputs of a downstream simulator. Finally, in sensitive applications, one must also assess privacy risks and memorization.

This last issue directly connects synthetic populations to current debates on synthetic data generated through deep learning or LLMs: good empirical fidelity does not imply freedom from information leakage, nor does it guarantee the absence of representational bias.

\section{Potential Contributions of LLMs}
\subsection{From profiles to behaviors}
LLMs were not originally designed for statistical population synthesis. Their main contribution lies instead in conditional text generation, including the production of preferences, arguments, explanations, or responses. This property has led several authors to use them as synthetic respondents or approximations of human sub-populations. \citet{argyle2023}, for instance, introduce the idea of \emph{silicon samples}, that is, artificial samples obtained by conditioning a model on sociodemographic descriptions. \citet{sun2024} show that aggregate information at the subgroup level may sometimes be sufficient to approximate certain opinion distributions.

The central result of this literature is not that LLMs generate representative populations in a strong statistical sense, but rather that they capture part of the regularities associated with demographic and discursive profiles embedded in their training data.

\subsection{LLMs and generative agents}
A second line of research uses LLMs not as survey-response generators, but as engines for generative agents. \citet{park2023} show that a system combining memory, reflection, and planning can produce apparently coherent behaviors within a small social environment. At a broader level, the survey by \citet{gao2024} synthesizes the emergence of LLM-enabled agent-based modeling and identifies four major challenges: environmental perception, human alignment, action generation, and evaluation.

For the problem of synthetic populations, the implication is straightforward: LLMs may be useful for endowing already calibrated agents with a more flexible, contextual, and linguistic behavioral layer than fixed rules can provide. However, they do not by themselves offer the guarantees of representativeness required for microstatistical reconstruction.

\subsection{Promising uses}
Three uses appear particularly promising:
\begin{enumerate}
  \item narrative or semantic imputation of hard-to-observe attributes, such as motivations, attitudes, or arguments;
  \item exploratory simulation of survey responses, especially in pretesting or hypothesis-generation settings \citep{sarstedt2024};
  \item behavioral animation of agents in models whose demographic structure has already been built through classical synthesis methods.
\end{enumerate}

In all three cases, LLMs serve an auxiliary role rather than a full substitute for population synthesis.

\section{Methodological Limits of LLMs for Synthetic Populations}
Several limitations call for strong caution. First, LLMs tend to produce smooth, normative, and prompt-sensitive responses. Within-group variance is often underestimated. Second, calibration to target distributions remains fragile: a demographic instruction does not guarantee compliance with margins, nor preservation of real correlations across variables. Moreover, biases present in training corpora may be reproduced or amplified, especially along dimensions such as gender, social background, or political culture \citep{argyle2023,gao2024}.

Another difficulty concerns reproducibility. Outputs depend on the model, its version, the temperature setting, the system prompt, and sometimes proprietary adjustments that are not documented. For demanding academic uses, this opacity is a serious obstacle. Finally, external validity remains limited: an LLM may appear credible in local interactions while still producing unrealistic behavioral distributions at the scale of an entire population.

\section{Toward Hybrid Architectures}
The state of the art suggests that hybrid architectures are currently the most defensible option. In such a scheme, population structure is generated through classical constrained methods: IPF, resampling, combinatorial optimization, or joint-distribution inference. LLMs then intervene to enrich latent attributes, produce conditional decisions, or simulate interactions among agents.

Such a division of roles offers three advantages. First, it preserves the statistical traceability of the population. Second, it allows richer behavioral heterogeneity than deterministic rules alone. Third, it facilitates auditing, since structural errors and behavioral errors can be assessed separately.

This perspective nevertheless implies a demanding research agenda: defining multi-level calibration protocols, coupling quantitative constraints with controlled prompting or fine-tuning strategies, and designing evaluation batteries capable of jointly assessing statistical fidelity and behavioral plausibility.

\section{Research Agenda}
Three questions appear especially important for future work.

The first is how to combine strong statistical constraints with LLM-based generation without losing higher-level coherence. This requires explicit mechanisms for constraint-aware generation.

The second is how to validate the latent dimensions produced by LLMs, for example opinions, motivations, or preferences. Direct comparisons with real surveys, activity diaries, or behavioral traces remain necessary.

The third is how to govern risks. Issues of privacy, bias, transparency, and dependence on proprietary models are central if these approaches are to support public decision-making or reproducible scientific results.

\section{Conclusion}
The literature on synthetic populations has reached methodological maturity through constraint-based calibration, multi-level coherence, and statistical evaluation. LLMs introduce a significant extension, but at another level: the production of behaviors, discourse, and conditional responses. The scientific challenge, therefore, is not to replace classical methods with LLMs, but to build hybrid pipelines in which statistical reconstruction guarantees a minimum degree of representativeness while LLMs enrich the dynamics of agents.

At the current state of the field, the most defensible use of LLMs is to complement an already calibrated synthetic population, not to substitute for the core synthesis process. This conclusion is methodologically conservative, but it is consistent with the available evidence: LLMs are promising for local realism and behavioral expressiveness, whereas classical approaches remain superior for statistical control, transparency, and auditability.

\bibliographystyle{plainnat}
\bibliography{references_populations_synthetiques_llms}

\end{document}
